% -------------------------------------------------------
%  Abstract
% -------------------------------------------------------

\شروع{وسط‌چین}
\مهم{چکیده}
\پایان{وسط‌چین}
\بدون‌تورفتگی

یکی از چالش‌های اصلی بازشناسی مقاوم گفتار فارسی، ناهمخوانی میان داده‌های آموزشی رایج و گفتار واقعی در تماس‌های تلفنی و سامانه‌های صدا روی پروتکل اینترنت است. این پایان‌نامه دو مسیر داده‌محور و بازتولیدپذیر را برای کاهش این ناهمخوانی و کمبود داده بررسی می‌کند: آموزش چندشرایطی با تخریب‌های مخابراتی و افزایش مقیاس آموزش با گفتار بلندِ شبه‌برچسب‌خورده. برای مسیر نخست، از گفتار فارسی جفت‌های هم‌تراز پاک--تخریب‌شده ساخته می‌شود و نویز محیطی، تغییر بهره، محدودیت پهنای‌باند، چرخهٔ واقعی کدک‌های صوتی از طریق \کد{ffmpeg} و افت بسته با بذرها و فرادادهٔ ثبت‌شده اعمال می‌شوند. سپس ویسپر-کوچک یک بار با داده‌های معمولی و گونه‌های بلند و بار دیگر با افزودن داده‌های تخریب‌شده ریزتنظیم می‌شود.

برای سنجش سود پیچیدگی معماری فراتر از آموزش چندشرایطی، یک مدل دوجریانه نیز طراحی می‌شود که بازنمایی‌های لاگ-مل اصلی و بهسازی‌شده را با رمزگذار مشترک ویسپر پردازش و با توجه متقاطع دوسویه و دروازهٔ زمان--ویژگی ترکیب می‌کند. بهساز این مدل یک U-Net باقیمانده با مدل‌سازی زمانی است و کل سامانه طی سه مرحله تا آموزش مشترک انتهابه‌انتها آماده می‌شود. در مسیر دوم، از میان \lr{3,777} کتاب صوتی ایران‌صدا با حدود \lr{6,300} ساعت صوت خام، \lr{250} ساعت برای قطعه‌بندی، شبه‌برچسب‌گذاری با ویسپر-متوسطِ سازگارشده با فارسی، اصلاح محافظه‌کارانه و پالایش انتخاب می‌شود. داده‌های پذیرفته‌شده بدون تغییر معماری به ترکیب کامل آموزش چندشرایطی ویسپر-کوچک افزوده می‌شوند. یک ویسپر-متوسط فارسی نیز به‌صورت آزمایشی جداگانه برای بررسی نتیجهٔ افزایش ظرفیت مدل ارزیابی می‌شود.

مدل‌ها روی پنج مجموعه‌دادهٔ فارسی، از جمله گفتار واقعی تلفنی و \lr{VoIP} در AGFarsdat، ارزیابی شدند. در مقایسهٔ کنترل‌شدهٔ ویسپر-کوچک، آموزش چندشرایطی نرخ خطای واژهٔ تجمیعی را از \lr{23.27\%} به \lr{17.83\%} و نرخ خطای نویسه را از \lr{10.84\%} به \lr{7.93\%} کاهش داد. مدل همجوشی با نرخ خطای واژهٔ \lr{18.44\%} از مدل معمولی بهتر بود، اما از مدل چندشرایطی عبور نکرد و نرخ خطای نویسهٔ آن به \lr{14.04\%} رسید. میانگین وزن دروازه برای نماهای بهسازی‌شده و اصلی به‌ترتیب \lr{0.0039} و \lr{0.9961} بود که اتکای تقریباً کامل مدل به جریان اصلی را نشان می‌دهد. هرچند همجوشی در AGFarsdat نسبت به مدل چندشرایطی \lr{0.42} واحد درصد نرخ خطای واژهٔ کمتری ثبت کرد، بدون آزمون آماری نمی‌توان این اختلاف کوچک را برتری پایدار دانست.

افزودن دادهٔ ایران‌صدا نرخ خطای واژه و نویسهٔ تجمیعی ویسپر-کوچک چندشرایطی را به‌ترتیب به \lr{17.22\%} و \lr{7.70\%} رساند، اما اثر آن حوزه‌وابسته بود: عملکرد در سه مجموعهٔ عمومی بهتر شد، درحالی‌که نرخ خطای واژه در AGFarsdat از \lr{25.80\%} به \lr{26.66\%} افزایش یافت. ویسپر-متوسط با نرخ خطای واژهٔ \lr{14.36\%} و نرخ خطای نویسهٔ \lr{6.68\%} بهترین نتیجهٔ تجمیعی مشاهده‌شده را به دست آورد، ولی به‌دلیل ظرفیت بیشتر، اختلاف آن با مدل‌های کوچک یک مقایسهٔ کنترل‌شدهٔ صرفاً داده‌محور نیست. در مجموع، داده‌افزایی مخابراتیِ قابل ردیابی و آموزش چندشرایطی روشن‌ترین خط مبنا برای مقاومت مخابراتی را فراهم کردند؛ معماری همجوشی در شکل فعلی سود تجمیعی بیشتری نداشت و گفتار بلند شبه‌برچسب‌خورده عمدتاً تعمیم در دامنه‌های عمومی را بهبود داد، نه گفتار مخابراتی واقعی را.

\پرش‌بلند
\بدون‌تورفتگی \مهم{کلیدواژه‌ها}:
بازشناسی خودکار گفتار فارسی، گفتار تلفنی، صدا روی پروتکل اینترنت، ویسپر، آموزش چندشرایطی، شبه‌برچسب‌گذاری، گفتار بلند، همجوشی دوجریانه
\صفحه‌جدید
